\section{Problem Formulation}\label{sec:problem}
The long-only mean-variance portfolio solves
\[
\min_{w}\; w^\top\Sigma\, w - \rho\,\mu^\top w
\quad\text{subject to}\quad
\mathbf{1}^\top w = 1,\quad w\geq 0~,
\]
where $\Sigma$ is the $n\times n$ covariance matrix of returns, $\mu$ is the $n\times 1$ vector of expected returns, and $\rho\geq 0$ controls the risk-return trade-off; setting $\rho=0$ recovers the pure long-only minimum variance portfolio.
Since both $\Sigma$ and
$\mu$ are unknown, in practice they must be replaced by estimates. Let $X\in\mathbb{R}^{T\times n}$ denote the demeaned return matrix for $n$ assets over the previous $T$ trading periods. The feasible problem formulation then is
\[
\min_{w}\; w^\top\hat\Sigma\, w - \rho\,\hat\mu^\top w
\quad\text{subject to}\quad
\mathbf{1}^\top w = 1,\quad w\geq 0~.
\]
The classical choice for estimating $\Sigma$ is
the sample covariance matrix $\hat\Sigma = X^\top X/T$, for which $w^\top\hat\Sigma\, w = \norm{Xw}^2/T$, and the problem becomes
\[
\min_{w}\;\frac{\norm{Xw}^2}{T} - \rho\,\hat\mu^\top w
\quad\text{subject to}\quad
\mathbf{1}^\top w = 1,\quad w\geq 0~.
\]
This reframing is key: the variance term is now a normalised squared norm of $Xw$, and the gradient $w\mapsto X^\top(Xw)$ can be evaluated with two matrix-vector products, without ever forming $X^\top X$.

\begin{table}[ht]
\centering
\begin{tabular}{lll}
\toprule
Parameter & Type & Description \\
\midrule
$X$        & $\mathbb{R}^{T \times n}$ & Demeaned return matrix \\
$\rho$     & $\mathbb{R}_{\geq 0}$     & Return-tilt weight (default $0$) \\
$\mu$      & $\mathbb{R}^n$            & Expected returns (default $\mathbf{0}$) \\
\bottomrule
\end{tabular}
\caption{Parameters of the \texttt{Problem} specification.}
\label{tab:params}
\end{table}

\paragraph{Dual problem}
Including multipliers for both constraints, the full Lagrangian is
\[
\mathcal{L}(w,\, \lambda,\, \nu)
= \frac{\lVert Xw\rVert^2}{T} - \rho\hat\mu^\top w
  - \lambda(\mathbf{1}^\top w - 1) - \nu^\top w\;,
\]
where $\lambda\in\mathbb{R}$ is the budget multiplier and $\nu\geq\mathbf{0}$ is the vector of
non-negativity multipliers. The stationarity condition
\[
\tfrac{2}{T}\, X^\top X w = \lambda\mathbf{1}+\nu+\rho\hat\mu
\]
gives $w^* = \tfrac{T}{2}(X^\top X)^{-1}s$ with $s = \lambda\mathbf{1}+\nu+\rho\hat\mu$.
Substituting into $\mathcal{L}$:
\[
\mathcal{L}(w^*,\,\lambda,\,\nu)
= \tfrac{1}{T}(w^*)^\top X^\top X w^* - s^\top w^* + \lambda\;.
\]
From stationarity $\tfrac{1}{T}X^\top X w^* = \tfrac{s}{2}$, so the first term equals
$\tfrac{1}{2}s^\top w^*$, giving
\[
g(\lambda,\,\nu)
= \tfrac{1}{2}s^\top w^* - s^\top w^* + \lambda
= -\tfrac{1}{2}s^\top w^* + \lambda\;.
\]
Substituting $w^* = \tfrac{T}{2}(X^\top X)^{-1}s$ yields the dual function
\[
g(\lambda,\,\nu)
= \inf_w\; \mathcal{L}(w,\,\lambda,\,\nu)
= \lambda - \tfrac{T}{4}\,
  (\lambda\mathbf{1}+\nu+\rho\hat\mu)^\top(X^\top X)^{-1}
  (\lambda\mathbf{1}+\nu+\rho\hat\mu)\;,
\]
defined when $X$ has full column rank; the dual problem is
$\max_{\lambda,\,\nu\geq\mathbf{0}}\; g(\lambda,\,\nu)$.
The feasible set $\{w:\mathbf{1}^\top w=1,\,w\geq\mathbf{0}\}$ is compact and
$w=\tfrac{1}{n}\mathbf{1}$ is strictly feasible, so Slater's condition holds and
strong duality obtains.

At the optimum, complementary slackness ($\nu_i w_i = 0$ for all $i$) partitions assets
into two types. \emph{Active} assets ($w_i>0$) have $\nu_i=0$; stationarity then gives
\[
\tfrac{2}{T}\,(X^\top Xw)_i = \lambda+\rho\hat\mu_i\;,
\]
equating each active asset's marginal variance to the budget shadow price adjusted for its
expected return. \emph{Excluded} assets ($w_i=0$) satisfy
\[
\nu_i = \tfrac{2}{T}\,(X^\top Xw)_i - \lambda - \rho\hat\mu_i \geq 0~.
\]
Their marginal variance exceeds this threshold, so holding them at a positive weight would
increase variance net of return. The primal-dual loop in Section~\ref{sec:primaldual} is a direct
implementation of these conditions: it terminates precisely when primal and dual feasibility
both hold, which (together with stationarity from each inner solve) is sufficient for
global optimality.
